\documentclass[11pt]{article}
\usepackage[margin=1.2in]{geometry}
\usepackage{amsmath,amssymb,amsthm}
\usepackage{xspace}
\usepackage[T1]{fontenc}
\usepackage{lmodern}
\usepackage{microtype}
\usepackage[hidelinks]{hyperref}

% --- Macros ---
\newcommand{\RetroDiff}{\textsc{RetroDiff}\xspace}
\newcommand{\Lcal}{\mathcal{L}}
\newcommand{\Ecal}{\mathbb{E}}
\newcommand{\pdata}{p_{\mathrm{data}}}
\newcommand{\ptheta}{p_\theta}
\newcommand{\ELBO}{\mathrm{ELBO}}
\newcommand{\KL}{\mathrm{KL}}

\pagestyle{empty}

\begin{document}

\noindent
The training objective of \RetroDiff derives from the variational lower bound (ELBO)
on the log-likelihood of clean data $x$.
Applying Jensen's inequality to $\log \ptheta(x)$ yields

\begin{align}
  \log \ptheta(x)
    &\;\geq\;
    \Ecal_{q(x_{1:T}\mid x)}\!\left[
      \log\frac{\ptheta(x_{0:T})}{q(x_{1:T}\mid x)}
    \right]
    = -\Lcal_{\ELBO},
  \label{eq:elbo}
\end{align}

\noindent
where $q$ is the fixed forward noising process that adds Gaussian noise over $T$ steps,
and $\ptheta$ is the learned reverse process.
The bound is tight when $q(\,\cdot\mid x)=\ptheta(\,\cdot\mid x)$;
in practice we minimise $\Lcal_{\ELBO}$ as a proxy for maximising $\log\ptheta(x)$.

Markov-decomposing $\ptheta(x_{0:T})$ and $q(x_{1:T}\mid x)$ step by step,
and collecting terms by timestep, converts \eqref{eq:elbo} into
a reconstruction term plus a sum of per-step KL divergences:

\begin{align}
  \Lcal_{\ELBO}
    &= \underbrace{\Ecal_q\!\bigl[-\log \ptheta(x_0\mid x_1)\bigr]}_{\text{reconstruction}}
     + \sum_{t=2}^{T}
       \underbrace{
         \KL\!\bigl(
           q(x_{t-1}\mid x_t,\,x)
           \;\|\;
           \ptheta(x_{t-1}\mid x_t)
         \bigr)
       }_{\text{step-}t\text{ denoising}}.
  \label{eq:decomp}
\end{align}

\noindent
Each KL term penalises the reverse model $\ptheta(x_{t-1}\mid x_t)$
for deviating from the analytic posterior
$q(x_{t-1}\mid x_t,x)$, which is Gaussian with closed-form mean
when the forward process is linear-Gaussian.
Equation~\eqref{eq:decomp} is standard for denoising diffusion probabilistic models;
the only free parameters are the neural-network weights $\theta$.

\RetroDiff's contribution is to condition each denoising step on a set of
$k$ exemplars $R(x)$ retrieved from the training corpus given $x$.
Averaging \eqref{eq:decomp} over the data distribution
with the retrieval conditioning in place gives the \RetroDiff objective:

\begin{align}
  \Lcal_{\RetroDiff}(\theta)
    &= \Ecal_{x\sim\pdata}\!\left[
        \Lcal_{\ELBO}\!\bigl(x \mid R(x)\bigr)
       \right].
  \label{eq:retro}
\end{align}

\noindent
Here $R(x)$ denotes the $k$-nearest-neighbour exemplar set retrieved for sample $x$,
and the ELBO in \eqref{eq:retro} is computed with $\ptheta(\,\cdot\,\mid\,\cdot,\,R(x))$
replacing the unconditional reverse process.
Minimising \eqref{eq:retro} encourages the model to exploit retrieved structure
while retaining the well-understood ELBO convergence guarantee of
\eqref{eq:elbo}--\eqref{eq:decomp}. \hfill$\square$

\end{document}
